% Test points -- nodes identified from the schematic. Expected signal
% levels are left for the simulation/measurement results.
\begin{table}[!htbp]
\centering
\caption{Designated test points.}
\label{tab:test-points}
\begin{tabularx}{\textwidth}{@{}lXX@{}}
\toprule
\textbf{Test point} & \textbf{Node (as drawn)} & \textbf{Observed on} \\
\midrule
TP1 & U1A non-inverting input (R1/R2 junction) & XSC1 channel A \\
TP2 & U1A output / U1B non-inverting input & XSC1 channel B \\
TP3 & U1B output (D1 cathode side) & --- \\
TP4 & C1/R4 hold node (D1 anode, U1B inverting input, U2A non-inverting input) & --- \\
TP5 & U2A output / Q1 gate & --- \\
\bottomrule
\end{tabularx}
\par\vspace{6pt}
\begin{minipage}{0.96\textwidth}
\raggedright\small \placeholder{Expected signal levels at each test point to be filled in from simulation and measurement; the XSC1 channel assignments should be confirmed against the Multisim file}
\end{minipage}
\end{table}
